\documentclass[11pt,a4paper]{article}
\usepackage[utf8]{inputenc}
\usepackage[T1]{fontenc}
\usepackage[spanish,es-nodecimaldot]{babel}
\usepackage{amsmath,amssymb}
\usepackage{graphicx}
\usepackage[margin=2.6cm]{geometry}

\title{El valle de la CCF en $h=-4$ vive en las crisis}
\author{}
\date{}

\begin{document}
\maketitle
\vspace{-3em}

\section{La correlaci\'on cruzada}

Con las tasas de variaci\'on log trimestrales
$p_t = 100\,[\ln P_t - \ln P_{t-1}]$ (ganancias) e
$i_t = 100\,[\ln I_t - \ln I_{t-1}]$ (inversi\'on)\footnote{$P_t$ e $I_t$ en
d\'olares corrientes. El c\'alculo rehecho con series deflactadas da el mismo
valle ($\rho(-4)=-0{,}233$ vs $-0{,}230$ nominal). Adem\'as, un componente
inflacionario com\'un a ambas series empuja los productos de desv\'ios hacia
\emph{positivo} (en d\'ecadas de inflaci\'on alta las dos variaciones nominales
quedan por encima de sus medias a la vez, y en d\'ecadas de inflaci\'on baja,
por debajo a la vez), de modo que trabajar en nominales juega en contra del
valle negativo, no a favor. Como tercera verificaci\'on, recalculando las
contribuciones contra una media m\'ovil de 10 a\~nos (que elimina cualquier
deriva inflacionaria lenta) el reparto es el mismo: 88\,\% de la covarianza
negativa en las vecindades de crisis, contra 87\,\% con la media global.}, la funci\'on de
correlaci\'on cruzada es
\[
\rho(h) = \operatorname{corr}\bigl(p_t,\, i_{t+h}\bigr).
\]
En $h=-4$ el par es \emph{(ganancia de hoy, inversi\'on de hace cuatro
trimestres)}.

Numeramos los trimestres de la muestra $t = 1, 2, \dots, 312$, donde $t=1$ es
1948Q2 y $t=312$ es 2026Q1. El par $(p_t,\, i_{t-4})$ solo puede formarse
cuando $t-4 \geq 1$, es decir para $t = 5, 6, \dots, 312$: los primeros cuatro
trimestres no tienen con qui\'en emparejarse y quedan $312-4=308$ pares.
$\rho(-4)$ es la correlaci\'on de Pearson com\'un sobre esos 308 pares:
\begin{equation}
\rho(-4)
= \frac{\displaystyle\sum_{t=5}^{312} \bigl(p_t - \bar p\bigr)\bigl(i_{t-4} - \bar\imath\bigr)}
       {\sqrt{\displaystyle\sum_{t=5}^{312} \bigl(p_t - \bar p\bigr)^2}\;
        \sqrt{\displaystyle\sum_{t=5}^{312} \bigl(i_{t-4} - \bar\imath\bigr)^2}}
\label{eq:ccf}
\end{equation}
con $\bar p$ el promedio de los $p_t$ usados ($t = 5,\dots,312$) y
$\bar\imath$ el de los $i_{t-4}$ usados ($t-4 = 1,\dots,308$).

El numerador no suma tasas: suma \emph{productos de desv\'ios}, un ``voto'' con
signo por trimestre --- negativo cuando inversi\'on alta precede ganancia baja
(o viceversa), positivo cuando ambas se desv\'ian para el mismo lado. El
denominador normaliza para que el cociente quede en $[-1,1]$. Sobre los 308
pares, $\rho(-4) = -0{,}23$.

\section{El experimento: ¿d\'onde se genera el $-0{,}23$?}

Cada trimestre tiene su \emph{contribuci\'on} $c_j$ --- un punto de la l\'inea
roja de la figura~\ref{fig:contrib} --- y la CCF es el promedio de las 308
contribuciones. La cuenta para un trimestre gen\'erico $j$ (con
$5 \leq j \leq 312$) es:
\begin{enumerate}
\item Tomar la tasa de variaci\'on de la ganancia del trimestre $j$: $p_j$.
\item Estandarizarla con la media y el desv\'io de \emph{toda} la serie de
variaciones de la ganancia, calculados una sola vez sobre los 312 trimestres:
$z^p_j = (p_j - \bar p)/s_p$.
\item Tomar la tasa de variaci\'on de la inversi\'on de cuatro trimestres
antes: $i_{j-4}$.
\item Estandarizarla con la media y el desv\'io de \emph{toda} la serie de
variaciones de la inversi\'on: $z^i_{j-4} = (i_{j-4} - \bar\imath)/s_i$.
\item Multiplicar:
\[
c_j = z^p_j \cdot z^i_{j-4}.
\]
\end{enumerate}
El signo de $c_j$ es el ``voto'' del trimestre: negativo si la inversi\'on
estaba por encima de su media y la ganancia cuatro trimestres despu\'es qued\'o
por debajo (o viceversa); positivo si ambas se desviaron para el mismo lado.

En s\'imbolos, la cuenta completa para un $j$ gen\'erico es
\[
z^p_j = \frac{p_j - \bar p}{s_p},
\qquad
z^i_{j-4} = \frac{i_{j-4} - \bar\imath}{s_i},
\qquad
c_j = z^p_j \cdot z^i_{j-4},
\]
donde el promedio y el desv\'io est\'andar de cada serie se calculan una sola
vez, sobre los 312 trimestres de la muestra:
\[
\bar p = \frac{1}{312}\sum_{t=1}^{312} p_t,
\qquad
s_p = \sqrt{\frac{1}{312}\sum_{t=1}^{312} \bigl(p_t - \bar p\bigr)^2},
\]
\[
\bar\imath = \frac{1}{312}\sum_{t=1}^{312} i_t,
\qquad
s_i = \sqrt{\frac{1}{312}\sum_{t=1}^{312} \bigl(i_t - \bar\imath\bigr)^2}.
\]
Ejemplo con $j = 100$, que es 1973Q1 (y $j-4 = 96$ es 1972Q1): es la misma
cuenta con los n\'umeros puestos,
\[
z^p_{100} = \frac{p_{100} - \bar p}{s_p} = \frac{11{,}21 - 1{,}55}{6{,}86} = +1{,}41,
\qquad
z^i_{96} = \frac{i_{96} - \bar\imath}{s_i} = \frac{7{,}08 - 1{,}53}{4{,}84} = +1{,}15,
\]
\[
c_{100} = z^p_{100} \cdot z^i_{96} = 1{,}41 \times 1{,}15 = +1{,}62 :
\]
pleno boom previo a la crisis de 1975, ganancia e inversi\'on rezagada ambas
por encima de sus medias, voto positivo. En un trimestre de crisis pasa lo
contrario: la inversi\'on de $j-4$ (fin de la expansi\'on) est\'a por encima de
su media y la ganancia de $j$ por debajo, y el producto sale negativo.

\begin{figure}[t]
\centering
\includegraphics[width=0.8\textwidth]{fig_ccf_movil.pdf}
\caption{Arriba: las tasas $p_t$ e $i_t$; el sombreado marca las recesiones
NBER. Abajo: la contribuci\'on $c_j$ de cada trimestre, coloreada seg\'un el
caso. En gris las positivas (ambas variaciones del mismo lado de su media). En
rojo las negativas de \emph{sobreacumulaci\'on}: inversi\'on por encima de su
media cuatro trimestres antes, ganancia por debajo hoy. En violeta las
negativas de \emph{recuperaci\'on}: inversi\'on derrumbada cuatro trimestres
antes, ganancia rebotando hoy (como 1950 o 2021). Las rojas se apilan entrando
a las crisis y las violetas saliendo: son las dos mitades del mismo ciclo
presa--depredador, vistas con medio ciclo de distancia.}
\label{fig:contrib}
\end{figure}

Repitiendo la cuenta para los 308 trimestres y promediando se recupera la CCF:
\[
\rho(-4) = \frac{1}{308}\sum_{j=5}^{312} c_j = -0{,}224 \approx -0{,}23
\]
(la peque\~na diferencia con \eqref{eq:ccf} es que all\'i las medias se
calculan sobre los 308 pares y aqu\'i sobre los 312 trimestres; num\'ericamente
despreciable). El experimento consiste en separar los 308 votos en dos urnas
antes de sumar:
\begin{enumerate}
\item Para cada una de las 12 recesiones NBER se ubica su \emph{fondo}: el
trimestre con la tasa de inversi\'on m\'as baja dentro de la ventana que va del
pico NBER menos 2 trimestres al valle NBER m\'as 4. Los fondos resultantes son
1949Q2, 1953Q4, 1957Q4, 1960Q4, 1970Q4, 1975Q1, 1980Q2, 1982Q1, 1990Q4,
2001Q4, 2009Q1 y 2020Q2.
\item Se marca como ``vecindad de crisis'' todo trimestre a distancia de a lo
sumo 6 del fondo m\'as cercano (hasta 13 trimestres por crisis; las ventanas de
1958--1961 y de 1980--1982 se solapan y cada trimestre se cuenta una sola vez).
De los 308 trimestres con par, quedan marcados 143.
\item Se suman los $c_j$ de los trimestres marcados por un lado y los del resto
por el otro.
\end{enumerate}
El reparto es:
\[
\underbrace{\sum_{j=5}^{312} c_j = -69}_{\text{los 308 trimestres}}
\;=\;
\underbrace{-60}_{\substack{\text{vecindades de crisis}\\ 143 \text{ trim}}}
\;+\;
\underbrace{-9}_{\substack{\text{resto de la muestra}\\ 165 \text{ trim}}}
\]
El 87\% de la covarianza negativa se genera en las vecindades de las crisis.
El mismo resultado en t\'erminos de correlaci\'on por r\'egimen: aplicando
\eqref{eq:ccf} solo a los pares cuyos dos trimestres ($t$ y $t-4$) caen dentro
de un mismo bloque de vecindad, $\rho(-4)=-0{,}28$ ($-0{,}38$ excluyendo 2008 y
2020); con los pares enteramente fuera de las vecindades, $0{,}00$.

\noindent\emph{Salvedad.} La correlaci\'on no discrimina mecanismos por s\'i
sola; ese trabajo lo hacen los controles (la regresi\'on que condiciona por la
historia propia de la ganancia y el contraste con series sustitutas sin
acople).

\end{document}
